% DEFINED:   sec:link-likelihood, sec:mdn
% RESOLVED:  sec:autograd-note
% FORWARD:   (none)

\section{Link and likelihood are independent choices}\label{sec:link-likelihood}

The catalogue in \Cref{sec:catalogue} paired each likelihood with its
canonical link. That is the \emph{natural} pairing, the one that makes
the gradient collapse to $\hat{p} - y$ and the optimization land
cleanly. It is not the \emph{only} pairing.

Suppose the target $y$ is a continuous proportion: $y \in [0, 1]$,
something like ``fraction of pixels lit,'' ``vote share,'' or ``relative
humidity.'' Not a class label, a \emph{number}. We want the prediction
$\hat{p} \in (0, 1)$ as well.

The natural link to enforce $\hat{p} \in (0, 1)$ is the sigmoid:
$\hat{p} = \sigma(z)$. We can pair sigmoid with two different
likelihoods:

\textbf{Sigmoid plus MSE on the sigmoid output.} Loss:
$L = \tfrac{1}{2}(y - \sigma(z))^2$. Implicit likelihood: Gaussian
noise in proportion space. Gradient: $(\sigma(z) - y) \cdot \sigma'(z)$,
the chain rule. This works fine in the bulk of the domain but saturates
near $z \to \pm\infty$ (where $\sigma'(z) \to 0$ and gradient flow
stalls), making it hard for the model to push predictions very near the
boundaries.

\textbf{Sigmoid plus Beta NLL.} Predict the mean
$\mu = \sigma(z_\mu)$ and a ``precision'' $\nu = \mathrm{softplus}(z_\nu)$,
set $\alpha = \mu\nu$, $\beta = (1 - \mu)\nu$, and use the Beta NLL,
\[
  L = -(\alpha - 1)\log y - (\beta - 1)\log(1 - y) + \log B(\alpha, \beta).
\]
This treats $y$ as drawn from a Beta distribution. It is far more
sensitive than the MSE choice near the boundaries (where Beta
log-density is heavy-tailed in the right way for proportion data) and
predicts a full posterior over $[0, 1]$ instead of a point.

Same link (sigmoid). Different likelihoods (Gaussian on the sigmoid
output versus Beta directly). The choice changes what is assumed about
the noise distribution, how the loss treats extreme values, and whether
the head is one output or two.

This breaks the catalogue's apparent one-to-one. The link function is
not uniquely determined by the assumed distribution. There is a canonical
pairing, the one the theorem gives, but it is one choice among many.
The skill is recognizing the choice.

Beta NLL is not in \texttt{scratchnn}, but it is a natural extension:
it requires only \texttt{math.lgamma}, which is in the standard library.
The components (sigmoid, softplus, and the Beta log-density) all follow
the same \texttt{Loss} interface as the losses already built. The
pedagogical point of this section is that link and likelihood are
independent, not that every likelihood needs to be in the library.

\section{Beyond unimodal: a pointer to mixture density networks}\label{sec:mdn}

All of the heads above assumed $p(y \mid x)$ is unimodal. When it is
not, a unimodal head collapses: a robot arm with two valid joint
configurations for the same end-effector pose, an inverse problem with
multiple solutions, the conditional next-frame distribution for a
stochastic video model. A Gaussian head will predict the mean of the
modes, which may lie in a region of \emph{zero density} under the true
distribution.

A \textbf{mixture density network (MDN)} outputs the parameters of a
$K$-component Gaussian mixture: $K$ means, $K$ scales, and $K$ mixture
weights (the latter through a softmax head). The NLL is
\[
  -\log \sum_{k=1}^{K} \pi_k(x)\,\mathcal{N}\!\bigl(y \mid \mu_k(x), \sigma_k(x)^2\bigr).
\]
Standard regression collapses to the midpoint between two modes; the
MDN captures each mode with its own component, weighted by its own
probability.

The MDN is not in \texttt{scratchnn}. It is a natural extension, framed
the same way \Cref{sec:autograd-note} frames automatic differentiation:
the idea is the same framework (output head plus matching NLL), the
implementation composes pieces the library already has (softplus for the
scales, softmax for the weights, logsumexp for the stable mixture
log-density), plus the hand-derived gradient through the mixture.
Building it is the next thing to build, not a new concept.
